% Options for packages loaded elsewhere
\PassOptionsToPackage{unicode}{hyperref}
\PassOptionsToPackage{hyphens}{url}
\documentclass[
  11pt,
]{article}
\usepackage{xcolor}
\usepackage[margin=25mm]{geometry}
\usepackage{amsmath,amssymb}
\setcounter{secnumdepth}{-\maxdimen} % remove section numbering
\usepackage{iftex}
\ifPDFTeX
  \usepackage[T1]{fontenc}
  \usepackage[utf8]{inputenc}
  \usepackage{textcomp} % provide euro and other symbols
\else % if luatex or xetex
  \usepackage{unicode-math} % this also loads fontspec
  \defaultfontfeatures{Scale=MatchLowercase}
  \defaultfontfeatures[\rmfamily]{Ligatures=TeX,Scale=1}
\fi
\usepackage{lmodern}
\ifPDFTeX\else
  % xetex/luatex font selection
\fi
% Use upquote if available, for straight quotes in verbatim environments
\IfFileExists{upquote.sty}{\usepackage{upquote}}{}
\IfFileExists{microtype.sty}{% use microtype if available
  \usepackage[]{microtype}
  \UseMicrotypeSet[protrusion]{basicmath} % disable protrusion for tt fonts
}{}
\makeatletter
\@ifundefined{KOMAClassName}{% if non-KOMA class
  \IfFileExists{parskip.sty}{%
    \usepackage{parskip}
  }{% else
    \setlength{\parindent}{0pt}
    \setlength{\parskip}{6pt plus 2pt minus 1pt}}
}{% if KOMA class
  \KOMAoptions{parskip=half}}
\makeatother
\setlength{\emergencystretch}{3em} % prevent overfull lines
\providecommand{\tightlist}{%
  \setlength{\itemsep}{0pt}\setlength{\parskip}{0pt}}
\usepackage{bookmark}
\IfFileExists{xurl.sty}{\usepackage{xurl}}{} % add URL line breaks if available
\urlstyle{same}
\hypersetup{
  hidelinks,
  pdfcreator={LaTeX via pandoc}}

\author{}
\date{}

\begin{document}

\section{Bott periodicity in KK: the Bott and Dirac
elements}\label{bott-periodicity-in-kk-the-bott-and-dirac-elements}

\emph{Self-checked by the writing AI, GPT-6.1 Sol (OpenAI). Public
domain (CC0).}

Bott periodicity comes from two concrete cycles. One multiplies by the
position vector in a Clifford algebra. The other differentiates forms on
Euclidean space. Their first product is a confining operator with one
even Gaussian zero mode. A rotation of two copies of Euclidean space
proves the reverse product. We keep the Clifford actions and the planar
projection explicit, so the orientation can be checked.

We use the graded tensor and Clifford Morita conventions of \emph{Graded
C*-algebras, Clifford algebras and graded Hilbert modules}, and the
product and associativity results of \emph{Homotopy, associativity, the
index pairing and KK-equivalence}. Section 0 proves the
bounded-transform and analytic product criterion locally. Lemma 4.0
below proves the oscillator domain, compact resolvent and even Gaussian
kernel. The Clifford matrix conventions are Theorem 5.1 and Corollary
5.2 of the preceding graded lesson; its Lemma 6.0a, Theorem 6.2 and
Proposition 6.3 prove the module tensor, stabilization and compact
tensor rules. Its Lemma 2.0 proves the compact ideal, multiplier
identification, countable-generation criterion and bounded strict
extension facts used by those module constructions and the technical
partition. The actual Morita product and its inverse tensor unitaries
are Proposition 7.1 of \emph{Connections and the existence of the
Kasparov product}. None of these inputs uses an unproved
Atiyah--Bott--Shapiro assertion.

\subsection{0. The analytic product criterion used
below}\label{the-analytic-product-criterion-used-below}

We prove the analytic recognition argument here. Lemma 0.0 proves the
required resolvent, continuous functional-calculus and tensor-extension
foundations locally. The bounded product, technical partition and
essential-replacement proofs are the preceding programme lessons
\emph{Kasparov's technical theorem} and \emph{Connections and the
existence of the Kasparov product}. The following argument also treats a
nonessential first representation.

\textbf{Lemma 0.0 (resolvents, continuous calculus and tensor
extension).} Let \(E\) be a Hilbert \(B\)-module, and let \(D=D^*\) be a
closed densely defined operator for which \(1+D^2\) has dense range.
Then \(D\pm i\) are bijective from \(\operatorname{Dom}D\) to \(E\),
with mutually adjoint inverse resolvents of norm at most one. There is a
nondegenerate continuous functional calculus
\(C_0(\mathbb R)\to\mathcal L(E)\); it extends to bounded continuous
functions, and \[
\|(D-z)^{-1}\|\leq |\operatorname{Im}z|^{-1}
\quad(\operatorname{Im}z\ne0).
\] If \(Y\) is any \(B\)-\(C\) correspondence, the closure of
\(D\otimes1\) on \(\operatorname{Dom}D\odot_B Y\) is self-adjoint
regular on \(E\otimes_B Y\), and its bounded continuous functional
calculus is \(f(D)\otimes1\). No separability or essentiality of the
left action on \(Y\) is assumed.

\textbf{Proof: resolvents.} Symmetry cancels the mixed terms in \[
\langle(D\pm i)x,(D\pm i)x\rangle
=\langle Dx,Dx\rangle+\langle x,x\rangle.
\] Thus these operators are bounded below by one. Their ranges are
closed: a convergent sequence of images makes the preimages Cauchy, and
closedness of \(D\) then gives a preimage of the limit. The range of
each contains the dense range of \(1+D^2\), since on
\(\operatorname{Dom}D^2\) this is the product \((D+i)(D-i)\) in either
order. Their ranges are therefore all of \(E\). Write
\(R_\pm=(D\pm i)^{-1}\). For \(a,b\in E\), symmetry gives \[
\langle R_+a,b\rangle
=\langle R_+a,(D-i)R_-b\rangle
=\langle a,R_-b\rangle.
\] Hence \(R_+^*=R_-\). The identities \(DR_\pm=1\mp iR_\pm\), first on
their ranges, give \[
R_+-R_-=-2iR_+R_-=-2iR_-R_+.
\] In particular they commute. Put \(U=1-2iR_+\). The preceding
identities give \(UU^*=U^*U=1\), and \(1-U=2iR_+\) has dense range
\(\operatorname{Dom}D\).

\textbf{Proof: functional calculus.} We first prove the continuous
calculus of a unitary locally. For a Laurent polynomial
\(p(z)=\sum_{j=-m}^m a_jz^j\), put \(M=\sup_{|z|=1}|p(z)|\). Given
\(\xi\in E\) and \(N>2m\), the continuous \(E\)-valued polynomial \[
w_N(z)=N^{-1/2}\sum_{k=0}^{N-1}z^kU^{-k}\xi
\] has
\(\int_{\mathbb T}\langle w_N,w_N\rangle\,d\theta/(2\pi)=\langle\xi,\xi\rangle\).
For \(m\leq l\leq N-1-m\), the coefficient of \(z^l\) in \(p(z)w_N(z)\)
is \(N^{-1/2}U^{-l}p(U)\xi\). Integrating the inner product cancels
distinct monomials, and every omitted coefficient contributes a positive
element. Since pointwise \(|p(z)|^2\leq M^2\), this gives the
positive-order inequality \[
\frac{N-2m}{N}\langle p(U)\xi,p(U)\xi\rangle
\leq \int_{\mathbb T}\langle p w_N,p w_N\rangle\,\frac{d\theta}{2\pi}
\leq M^2\langle\xi,\xi\rangle.
\] Take norms and then \(N\to\infty\). Thus
\(\|p(U)\|\leq\|p\|_\infty\), for every Laurent polynomial, on a Hilbert
module as well as a Hilbert space.

Laurent polynomials are uniformly dense in \(C(\mathbb T)\). Indeed the
kernels \(K_N(\theta)=N^{-1}|\sum_{k=0}^{N-1}e^{ik\theta}|^2\) are
positive with normalized integral one. The geometric-series formula
bounds their integral outside any fixed neighborhood of zero by a
constant times \(N^{-1}\). Uniform continuity then makes convolution
with \(K_N\) converge uniformly to a continuous circle function; each
convolution is a Laurent polynomial. The polynomial norm bound therefore
extends evaluation at \(U\) to a unital contractive *-homomorphism
\(C(\mathbb T)\to\mathcal L(E)\). Multiplication and adjoints pass
through uniform polynomial approximation. No earlier normal-calculus or
spectral-representation assertion is needed for this construction. The
map \[
u(t)=(t-i)/(t+i)
\] identifies \(\mathbb R\) with the circle minus \(1\). For
\(f\in C_0(\mathbb R)\), its transported function extends by zero at
\(1\); evaluate it at \(U\) to define \(f(D)\). This is a contractive
*-homomorphism. It is nondegenerate because the function
\(1-u(t)=2i/(t+i)\) has operator value \(1-U\), whose range is dense.

For completeness, its extension to \(C_b(\mathbb R)\) is also
constructive. Let \(e_n\in C_c(\mathbb R)\) be positive cutoffs tending
uniformly to one on compact sets, with \(0\leq e_n\leq1\). For bounded
continuous \(h\), the uniformly bounded operators \((he_n)(D)\) converge
on every vector \(f(D)\xi\), because \(he_nf\to hf\) uniformly. Such
vectors have dense span, so the operators converge on every vector to a
bounded operator \(h(D)\). The same argument for \(\overline h\) gives
its adjoint. Products and adjoints pass to these uniformly bounded
vector limits, proving the extended *-homomorphism. Its value does not
depend on the chosen cutoffs, by the same dense-span argument.

The identity \(D=i(1+U)(1-U)^{-1}\) holds on the dense domain
\(\operatorname{ran}(1-U)\). For nonreal \(z\), the bounded scalar
function \((t+i)/(t-z)\) therefore shows that the value of
\((t-z)^{-1}\) has range in \(\operatorname{Dom}D\) and is a two-sided
inverse of \(D-z\). Its norm is bounded by the scalar supremum
\(|\operatorname{Im}z|^{-1}\). Polynomial identities on \(U\) and their
uniform limits show that this calculus agrees with the two original
resolvents and with the usual products of their functions. In particular
the value of \((1+t^2)^{-1}\) is \((1+D^2)^{-1}=R_+R_-\).

We also need the converse Cayley construction. If \(V\) is unitary and
\(1-V\) has dense range, then it is injective: its kernel equals the
kernel of \(1-V^*\), which is orthogonal to that dense range. On
\(\operatorname{ran}(1-V)\) define \[
S((1-V)x)=i(1+V)x.
\] Expanding both inner products and using \(V^*V=1\) proves symmetry.
The operators \(S+i\) and \(S-i\) have bounded, everywhere-defined,
mutually adjoint inverses \[
Q_+=(1-V)/(2i),\qquad Q_-=(V^*-1)/(2i).
\] These formulas prove that \(S\) is closed, since the inverse of an
injective bounded operator is closed on its range. If
\(y\in\operatorname{Dom}S^*\), choose \(x\in\operatorname{Dom}S\) with
\((S-i)x=(S^*-i)y\). Then \(y-x\in\ker(S^*-i)\), which is orthogonal to
the surjective range of \(S+i\). Hence \(y=x\), proving \(S=S^*\). The
commuting product \(Q_+Q_-\) has range in \(\operatorname{Dom}S^2\) and
is the inverse of \(1+S^2\), by the two resolvent identities. Thus \(S\)
is regular. This proves the converse without a graph-complement theorem.

\textbf{Proof: tensor extension.} Lemma 6.0a of Lesson 05 supplies the
unital *-homomorphism \(T\mapsto T\otimes1\) on adjointable operators.
Thus \(V=U\otimes1\) is unitary. The range of \(1-V\) is dense: the
algebraic tensors with first vector in \(\operatorname{Dom}D\) are
dense, and every such vector is in that range. Apply the converse Cayley
construction to \(V\), obtaining a self-adjoint regular \(S\). The
formulas for \(Q_\pm\) show that \(S(x\otimes y)=Dx\otimes y\) on the
algebraic domain. This domain is a graph core: its image under \(S+i\)
contains all elementary tensors and is dense, while \(S+i\), with its
bounded inverse and \(S Q_+=1-iQ_+\), identifies its graph norm with an
equivalent complete norm on \(E\otimes_B Y\). Consequently \(S\) is
precisely the closure of the algebraic tensor operator. Continuous
circle calculus commutes with \(T\mapsto T\otimes1\). The
bounded-function extension does too, since uniformly bounded vector
convergence on \(E\) gives vector convergence on elementary tensors and
hence on their dense span. This proves the claimed calculus identity. If
\(D\) is odd, its grading sends \(U\) to \(U^*\), and the same formulas
give \(f(D)\mapsto f(-D)\) and oddness of the tensor extension.
\(\square\)

The free
\href{https://www.math.uni-bonn.de/people/lesch/dl/pap/2011-KaaLes-ArXiv-v2.pdf}{Kaad--Lesch
author preprint, Section 2.3 and Proposition 4.1} supplies the
regular-operator formulation used here. Its references are not imported
as proof premises: the resolvent, Cayley, functional-calculus and tensor
assertions needed below have just been proved.

\textbf{Lemma 0.0a (the Euclidean Fourier facts used here).} With
normalization
\(\widehat f(\xi)=(2\pi)^{-n/2}\int e^{-ix\cdot\xi}f(x)\,dx\), Fourier
transformation is unitary on \(L^2(\mathbb R^n)\). On Schwartz functions
it sends \(\partial_j\) to multiplication by \(i\xi_j\), and
multiplication by \(x_j\) to \(i\partial_{\xi_j}\). The functions
\(z^k\), \(k\in\mathbb Z\), are a complete orthonormal basis of
\(L^2(\mathbb T,d\theta/(2\pi))\).

\textbf{Proof.} Integration by parts and differentiation under the
integral give the two intertwining formulas. Iterating them proves that
Fourier transformation preserves Schwartz space. The one-dimensional
Gaussian integral is \(\sqrt{2\pi}\): square the integral of
\(e^{-x^2/2}\) and use polar coordinates. For \(g(x)=e^{-x^2/2}\), its
transform solves \(\widehat g'=-\xi\widehat g\), by the intertwining
formulas, and its value at zero is one. Tensoring and positive real
rescaling give
\(\widehat {e^{-\epsilon|x|^2/2}}=\epsilon^{-n/2}e^{-|\xi|^2/(2\epsilon)}\).
For a Schwartz \(f\), insert \(e^{-\epsilon|\xi|^2/2}\) into its inverse
Fourier integral and use Fubini. The result is its convolution with the
positive mass-one Gaussian
\((2\pi\epsilon)^{-n/2}e^{-|x|^2/(2\epsilon)}\). This converges to
\(f\), by splitting the convolution into a small neighborhood and its
Gaussian tail. Dominated convergence on the inverse integral therefore
proves Fourier inversion. Fubini and inversion give
\(\int\overline{\widehat f}\widehat h=\int\overline f h\). Schwartz
functions are dense in \(L^2\): truncate an \(L^2\) function,
approximate the truncation by simple functions and smooth compactly
supported approximations, and mollify; translation continuity gives the
last convergence. Thus the isometry extends to \(L^2\); inversion on the
dense Schwartz space makes its range dense and hence onto.

For the circle, the monomials are orthonormal by direct integration.
Their span is dense without a completeness assumption: the Fejér kernels
\(K_N(\theta)=N^{-1}|\sum_{j=0}^{N-1}e^{ij\theta}|^2\) are positive,
have normalized integral one, and have integral tending to zero outside
every neighborhood of zero, from the geometric-series formula.
Convolution with them converges uniformly to each continuous circle
function, by uniform continuity and this tail bound, and each
convolution is a trigonometric polynomial. Continuous functions are
dense in circle \(L^2\), by approximation of step functions with
continuous ramps. This proves completeness. \(\square\)

An unbounded module \((E,\phi,D)\) means a countably generated graded
Hilbert module, a graded representation, and an odd self-adjoint regular
operator. We require \((1+D^2)^{-1}\phi(a)\) compact for every source
element. A dense graded \(*\)-subalgebra preserves
\(\operatorname{Dom}D\) and has bounded adjointable graded commutators
with \(D\). The representation need not be essential.

\textbf{Lemma 0.1 (bounded transformation with local compactness).} For
an unbounded module, \(F_D=D(1+D^2)^{-1/2}\) is a Kasparov-cycle
operator.

\textbf{Proof.} Put \(q(t)=(1+t^2)^{-1}\). Cutoffs \(e_n\) on \((0,1]\),
zero near zero and eventually one on every compact subinterval, make
\(e_n(q)/q\) bounded. Hence \(e_n(q(D))\phi(a)\) is compact. For
\(f\in C_0(\mathbb R)\), the functions \(fe_n(q)\) converge uniformly to
\(f\); therefore \(f(D)\phi(a)\) is compact. Adjointing gives
compactness on the other side.

For a homogeneous smooth source element write \(\epsilon=(-1)^{|a|}\),
\(c=[D,\phi(a)]_{\rm gr}\), and \(R_z=(D-z)^{-1}\). Domain preservation,
followed by multiplication by inverse resolvents, gives \[
R_z\phi(a)-\epsilon\phi(a)R_{\epsilon z}
=-\epsilon R_zcR_{\epsilon z},\qquad z=\pm i\rho.
\tag{0.1}
\] For \(Q_\rho=D(D^2+\rho^2)^{-1}=(R_{i\rho}+R_{-i\rho})/2\), it
follows that \(\|[Q_\rho,\phi(a)]_{\rm gr}\|\leq\|c\|\rho^{-2}\). The
commutator becomes compact after source multiplication on either side,
since the corresponding end resolvent is locally compact.

Functional calculus gives the vector integral \[
F_D\xi=\frac2\pi\int_0^\infty
D(1+D^2+t^2)^{-1}\xi\,dt.
\tag{0.2}
\] Indeed its partial-integral scalar functions are bounded by one,
converge uniformly on compact subsets of the real line, and converge
pointwise to \(x/\sqrt{1+x^2}\). Apply them first to the dense set of
vectors \(f(D)\eta\), \(f\in C_0(\mathbb R)\), and then use the uniform
bound. The uncommuted integral need not converge in operator norm. Its
commutator does: (0.1) bounds its integrand by \(\|c\|/(1+t^2)\). Thus
both source-localized commutators of \(F_D\) are compact. For
homogeneous smooth \(a,b\), the identity \[
[F_D,\phi(ab)]_{\rm gr}
=[F_D,\phi(a)]_{\rm gr}\phi(b)
 +(-1)^{|a|}\phi(a)[F_D,\phi(b)]_{\rm gr}
\] makes the unlocalized commutator for \(ab\) compact. Products from
the dense smooth algebra linearly span a dense subset of the source:
approximate an algebra approximate identity and the given element by
smooth elements and multiply. Since the commutator is bounded by
\(2\|a\|\) on each parity subspace, density proves compactness for every
source element. Finally \(F_D\) is odd and self-adjoint, and
\(F_D^2-1=-(1+D^2)^{-1}\). These are all cycle conditions. \(\square\)

Write \(E=E_1\widehat\otimes_B E_2\), \(\phi=\phi_1\widehat\otimes1\),
and \(E_1^0=\overline{\phi_1(A)E_1}\). For a homogeneous vector \(x\),
\(T_x\eta=x\widehat\otimes\eta\) is the creation operator.

\textbf{Lemma 0.2 (the connection estimate).} Suppose a homogeneous
linear subspace \(\mathcal X\subset\operatorname{Dom}D_1\), dense in
\(E_1^0\), satisfies the creation and adjoint domain inclusions and has
bounded adjointable errors \[
DT_x-(-1)^{|x|}T_xD_2,\qquad
T_x^*D-(-1)^{|x|}D_2T_x^*.
\tag{0.3}
\] Then \(F_D\) satisfies both \(F_{D_2}\)-connection conditions for
every \(x\in E_1^0\).

\textbf{Proof.} On \(E\oplus E_2\), the two creations are the
off-diagonal blocks of a homogeneous adjointable operator. Their domain
inclusions and (0.3) give a bounded graded commutator with
\(D\oplus D_2\). The inverse-resolvent identity (0.1) applies to this
operator too. Multiply its bounded-transform commutator on the right by
\(\operatorname{diag}(0,\phi_2(b))\), for a homogeneous smooth
second-source element \(b\). The localized end resolvent is compact, and
the bound \(C/(1+t^2)\) in (0.2) is integrable. The top-right block
proves \[
\bigl(F_DT_x-(-1)^{|x|}T_xF_{D_2}\bigr)\phi_2(b)
\] compact. Since \(T_{xb}=T_x\phi_2(b)\), moving \(F_{D_2}\) past
\(\phi_2(b)\) adds only the compact commutator proved in Lemma 0.1. Thus
the creation error for \(xb\) is compact; adjointing gives its second
condition. The linear span of these vectors is dense in \(E_1^0\), by
nondegeneracy of its right coefficient action, density of
\(\mathcal X\), and density of the smooth coefficient algebra. The
inequality \(\|T_x\|\leq\|x\|\) extends both assertions to every vector
of \(E_1^0\). \(\square\)

\textbf{Lemma 0.3 (bounded alignment, including essential connections).}
Let \(A\) be separable and the intermediate algebra be
\(\sigma\)-unital. Suppose \(F_1\) is a normalized first-cycle operator
and \(G\) is a self-adjoint contraction cycle on \(E\), with the
\(F_2\)-connection conditions on \(E_1^0\). If, for a fixed
\(0\leq\kappa<2\), \[
\phi(a)[F_1\widehat\otimes1,G]_{\rm gr}\phi(a)^*
\geq-\kappa\phi(aa^*)\pmod{\mathcal K(E)}
\qquad(a\in A),
\tag{0.4}
\] then \(G\) represents the Kasparov product.

\textbf{Proof.} First assume a full connection. The product-existence
proof, using this connection and its technical partition \(M+N=1\),
gives a product operator \(H=M^{1/4}SM^{1/4}+N^{1/4}GN^{1/4}\),
\(S=F_1\widehat\otimes1\). The same root commutation and defect
annihilation used there give, after a source sandwich, \[
[G,H]_{\rm gr}
\equiv M^{1/2}[G,S]_{\rm gr}+2N^{1/2}
\pmod{\text{source-localized compacts}}.
\] Consequently (0.4), \(0\leq M^{1/2}\leq1\), and positivity of the
second term bound this sandwich below by \(-\kappa\phi(aa^*)\).

Here are the quotient and homotopy details of that conclusion. Let
\(\mathcal D_\phi\) be the graded algebra of bounded operators whose
graded source commutators are compact, and let \(\mathcal C_\phi\) be
its ideal of two-sided source-locally compact operators. Closure under
products and adjoints follows from the graded Leibniz identity. For an
even self-adjoint \(Z\in\mathcal D_\phi\), positivity of every source
sandwich modulo compacts implies \(Z_-\in\mathcal C_\phi\). To see this,
work modulo \(\mathcal K(E)\), where \(Z\) commutes with the represented
source. The negative part of \(\phi(aa^*)Z\) is \(\phi(aa^*)Z_-\); it is
zero. Hence \(Z_-\phi(a)\) has zero product with its adjoint in this
quotient and is compact. Adjointing gives the other side. Continuous
functional calculus is legitimate because it follows by polynomial
approximation in the unital algebra generated by \(Z\).

Apply this to \(Z=[G,H]_{\rm gr}+\kappa\). In
\(\mathcal D_\phi/\mathcal C_\phi\), the cycle operators \(G,H\) are
self-adjoint unitaries and their anticommutator is at least \(-\kappa\).
For \(L_t=(1-t)G+tH\) this gives \[
L_t^2\geq (1-t)^2+t^2-\kappa t(1-t)
\geq(2-\kappa)/4.
\] Choose \(0<\delta<(2-\kappa)/4\) and normalize by
\(L_t\bigl(\max(L_t^2,\delta)\bigr)^{-1/2}\). This is a norm-continuous
path of cycles: it is an odd self-adjoint unitary in the defect
quotient, and its source commutators are compact by functional calculus.
At its endpoints normalization changes the original cycle only by a
two-sided locally compact operator. The straight interpolation for such
a change is a cycle, as follows by expanding its square and source
commutator. Thus \(G\) and \(H\) have the same class.

We now remove the full-connection assumption. Put
\(E^0=\overline{\phi(A)E}\); the balancing rule identifies it with
\(E_1^0\widehat\otimes_B E_2\). Proposition 5.1 of the product lesson
supplies essential replacement cycles \(F_1^0,G^0\) on these modules in
the original classes. Its connection relations for multiplication maps
\(A_a^1:E_1\to E_1^0\), \(A_a:E\to E^0\) are \[
F_1^0A_a^1-(-1)^{|a|}A_a^1F_1\in\mathcal K(E_1,E_1^0),\qquad
G^0A_a-(-1)^{|a|}A_aG\in\mathcal K(E,E^0).
\tag{0.5}
\] These maps are adjointable: their adjoints are multiplication by
\(a^*\) followed by inclusion. Normalizing the replacement cycles
preserves (0.5).

The operator \(G^0\) is a full connection on \(E_1^0\). Indeed
\(T_{ax}=A_aT_x\). Expanding its error with (0.5) leaves a compact term
and the source-localized creation error of \(G\). That latter term is
compact after inclusion in \(E\), because the connection error for
\(ax\in E_1^0\) is compact and the source commutator of \(G\) is
compact. It is compact with target \(E^0\) as well: for an adjointable
map \(R\), compactness of \(R^*R\) implies compactness of \(R\). In
fact, with \(q_\varepsilon(t)=t/(t+\varepsilon)\), the maps
\(Rq_\varepsilon(R^*R)\) are compact and converge in norm to \(R\),
since \(\sup_{t\geq0}\sqrt t\,\varepsilon/(t+\varepsilon)\to0\). Taking
adjoints and using density of the vectors \(ax\) proves both full
connection conditions.

For \(k\in\mathcal K(E_1)\), the weak connection also gives \[
\phi(a)[G,k\widehat\otimes1]_{\rm gr}\phi(b)\in\mathcal K(E).
\tag{0.6}
\] Indeed \(\phi_1(a)k\phi_1(b)\) is a norm limit of rank-one operators
with both vectors in \(E_1^0\). Their tensor images are \(T_xT_y^*\);
the connection identities cancel the two \(F_2\) terms in their
commutator. Expanding the commutator with the two outer source factors
adds only compact source commutators of \(G\). This proves (0.6).

Take even \(c\in A\). By (0.5),
\(R_c=(A_c^1)^*F_1^0A_c^1-\phi_1(c)^*F_1\phi_1(c)\) is compact. Put
\(S^0=F_1^0\widehat\otimes1\). Moving \(G^0\) past \(A_c\), then
sandwiching by \(\phi(d)\), gives \[
\phi(d)^*A_c^*[S^0,G^0]_{\rm gr}A_c\phi(d)
\equiv\phi(cd)^*[S,G]_{\rm gr}\phi(cd)
\pmod{\mathcal K(E)}.
\tag{0.7}
\] The extra bracket of \(R_c\widehat\otimes1\) is compact by (0.6); the
remaining errors are (0.5) and compact source commutators. Equation
(0.4), applied with \(a=(cd)^*\), transfers its lower bound through
(0.7).

This transfer holds on \(E^0\), not only after inclusion into \(E\). A
compact operator on \(E\) whose range and adjoint range lie in \(E^0\)
belongs to \(\mathcal K(E^0)\). For a proof, first regard it as an
adjointable map \(E\to E^0\); the compactness test just proved makes
that map compact. An approximate identity \(u_\lambda\) of
\(\mathcal K(E^0)\) then has \(u_\lambda T\to T\) in norm, by finite
rank approximation with target vectors in \(E^0\). Apply the same
argument to \(T^*\) to obtain \(Tu_\lambda\to T\). The sandwiches
\(u_\lambda T u_\lambda\) lie in \(\mathcal K(E^0)\) and converge to
\(T\). For the source sandwich in (0.7), its negative part has the same
range property, since its continuous defining function vanishes at zero.
Thus its compact negative part restricts compactly to \(E^0\). Positive
even approximate identities \(c\) give \(cd\to d\), proving (0.4) for
all sources on the essential module. The full-connection case applies
there. Essential replacement and homotopy invariance of products return
the assertion to the original modules. \(\square\)

\textbf{Lemma 0.4 (integrated positivity).} Let \(D,S\) be odd
self-adjoint regular operators with
\(\operatorname{Dom}D\subset\operatorname{Dom}S\). If \[
Q(\xi):=2\operatorname{Re}\langle D\xi,S\xi\rangle
\geq-c\langle\xi,\xi\rangle\quad(\xi\in\operatorname{Dom}D),
\tag{0.8}
\] for \(c\geq0\), then \([F_D,F_S]_{\rm gr}\geq-c\). Replacing \(S\) by
\(\alpha S\), \(\alpha>0\), replaces this bound by \(-\alpha c\).

\textbf{Proof.} For \(\lambda=1+u^2\), \(\mu=1+v^2\), set \[
\begin{gathered}
r_D=(D^2+\lambda)^{-1},\quad h_D=Dr_D,\quad k_D=\sqrt\lambda r_D,\\
r_S=(S^2+\mu)^{-1},\quad h_S=Sr_S,\quad k_S=\sqrt\mu r_S.
\end{gathered}
\] For \(\eta\in\operatorname{Dom}S\), expanding \(Dh_D=1-\lambda r_D\),
\(Dk_D=\sqrt\lambda h_D\) gives \[
Q(h_D\eta)+Q(k_D\eta)=2\operatorname{Re}\langle S\eta,h_D\eta\rangle.
\] The cancelling difference is
\(2\lambda\operatorname{Re}(\langle h_D\eta,Sr_D\eta\rangle-\langle r_D\eta,Sh_D\eta\rangle)\),
zero by symmetry of \(S\). All these vectors are in the indicated
domains by the domain inclusion. Apply the identity to
\(h_S\psi,k_S\psi\). Using \(Sh_S=1-\mu r_S\), \(Sk_S=\sqrt\mu h_S\),
the remaining two terms cancel by self-adjointness of \(h_D\). Hence \[
\sum_{w\in\{h_Dh_S,k_Dh_S,h_Dk_S,k_Dk_S\}}Q(w\psi)
=2\operatorname{Re}\langle\psi,h_Dh_S\psi\rangle.
\tag{0.9}
\] Integrate first over finite rectangles \(0\leq u,v\leq R\), with
factor \(4/\pi^2\). Formula (0.2) makes the right side converge to the
quadratic form of the bounded anticommutator. The lower bound (0.8)
bounds the left side below by minus \(c\) times the sum of the four
positive forms \(\langle w\psi,w\psi\rangle\). Their total integral is
at most \(\langle\psi,\psi\rangle\): first use \[
h_D^2+k_D^2=r_D,\qquad
\frac2\pi\int_0^\infty r_D\,du=(1+D^2)^{-1/2}\leq1,
\] then bound the two remaining sandwiches by \(h_S^2+k_S^2=r_S\) and
integrate once more. These increasing positive integrals also justify
passage to the rectangle limit without asserting absolute integrability
of the unsplit anticommutator integrand. The inequality holds for every
\(\psi\), hence as an operator inequality. Scaling \(S\) scales (0.8) by
\(\alpha\), proving the last assertion. \(\square\)

\textbf{Theorem 0.5 (the unbounded product criterion).} Let \(A\) be
separable and \(B\) be \(\sigma\)-unital. Suppose \((E_1,\phi_1,D_1)\),
\((E_2,\phi_2,D_2)\), and
\((E_1\widehat\otimes_B E_2,\phi_1\widehat\otimes1,D)\) are unbounded
modules. Suppose (0.3) holds on a homogeneous subspace of
\(\operatorname{Dom}D_1\) dense in \(\overline{\phi_1(A)E_1}\), with
both creation domain inclusions. Suppose
\(\operatorname{Dom}D\subset\operatorname{Dom}(D_1\widehat\otimes1)\),
and (0.8) holds with \(S=D_1\widehat\otimes1\). Then the
bounded-transform class of \(D\) is the Kasparov product of the classes
of \(D_1,D_2\). The first representation may be nonessential.

\textbf{Proof.} Lemma 0.1 supplies the cycles and Lemma 0.2 their
essential-source connection. Choose \(\alpha>0\) with \(\alpha c<2\).
Lemma 0.4 gives (0.4) for \(F_D\) and
\(F_{\alpha D_1}\widehat\otimes1\); functional calculus commutes with
the tensor representation. Lemma 0.3 identifies their product. Finally
positive rescaling preserves the class of the first operator: on a
compact positive interval of parameters, \(tx/\sqrt{1+t^2x^2}\) varies
uniformly in \(x\), continuously in \(t\). Its transforms form a
norm-continuous path of cycles by Lemma 0.1. Domains and bounded
commutators are unchanged by this positive scaling, and local
compactness follows from the vanishing-functions argument in its proof.
Thus \([\alpha D_1]=[D_1]\), proving the assertion. \(\square\)

The freely readable \href{https://arxiv.org/pdf/2006.10616v2}{van den
Dungen preprint, Sections 2.4 and 3.1}, printed pp.~11--16, provides the
comparison account for connection estimates and integrated positivity.
The proof used here is the complete local argument above. In particular
its nonessential representation step does not rely on the preprint's
reference to an external original proof.

\subsection{1. Two Clifford actions on
forms}\label{two-clifford-actions-on-forms}

Let \(C_n\) be the complex graded Clifford algebra on odd self-adjoint
generators \(e_j\), with \[
e_je_k+e_ke_j=2\delta_{jk}.
\] Put \[
B_n=C_0(\mathbb R^n)\widehat\otimes C_n.
\] The function algebra is trivially graded. On
\(\Lambda=\Lambda^*\mathbb C^n\), let \(\varepsilon_j\) be exterior
multiplication by the \(j\)-th basis vector and let
\(\iota_j=\varepsilon_j^*\). Grade forms by their degree modulo two.
Define \[
c_j=\varepsilon_j+\iota_j,\qquad
\widehat c_j=\varepsilon_j-\iota_j.
\tag{1.1}
\] The exterior-algebra relations give \[
\begin{gathered}
c_j^*=c_j,\quad \widehat c_j^*=-\widehat c_j,\\
\{c_j,c_k\}=2\delta_{jk},\quad
\{\widehat c_j,\widehat c_k\}=-2\delta_{jk},\quad
\{c_j,\widehat c_k\}=0 .
\end{gathered}
\tag{1.2}
\] Here braces are anticommutators. For example the mixed relation is
the sum of \(\{\varepsilon_j,\varepsilon_k\}\),
\(-\{\iota_j,\iota_k\}\), and the two cancelling Kronecker terms. All
operators in (1.1) are odd.

Thus \(e_j\mapsto c_j\) is a graded representation of \(C_n\). The
second family is deliberately skew-adjoint. Combining it with the
skew-adjoint derivative makes a self-adjoint Dirac operator.

\subsection{2. Multiplication by
position}\label{multiplication-by-position}

On the standard right module \(B_n\), multiplication on the left by \[
X(x)=\sum_{j=1}^n x_je_j
\] defines an odd self-adjoint regular operator. Its domain consists of
the \(b\in B_n\) for which \(Xb\in B_n\), and \[
(X\pm i)^{-1}(x)=\frac{X(x)\mp i}{1+|x|^2}.
\tag{2.1}
\] The inverse resolvents are bounded adjointable multipliers with dense
range: compactly supported Clifford-valued functions lie in their range.
Symmetry and the two mutually adjoint resolvents prove self-adjointness
and regularity by the converse resolvent/Cayley argument of local Lemma
0.0.

\textbf{Proposition 2.1 (the Bott cycle).} With the scalar source
representation, the bounded transform \[
F_n(x)=\frac{\sum_jx_je_j}{\sqrt{1+|x|^2}}
\tag{2.2}
\] gives a class \[
\beta_n=[(B_n,1,F_n)]\in KK(\mathbb C,B_n).
\]

\textbf{Proof.} The compact operators on the standard module \(B_n\) are
left multiplication by elements of \(B_n\). This follows on rank-one
operators from \(\theta_{b,d}(a)=bd^*a\), and the products \(bd^*\) span
a dense subspace. The operator in (2.2) is odd and self-adjoint,
commutes with the scalar action, and \[
F_n^2-1=-\frac1{1+|x|^2}.
\tag{2.3}
\] The right side belongs to \(B_n\), hence is compact. The module is
countably generated, since \(B_n\) is separable and has a countable
approximate identity. These are all cycle conditions. Alternatively
(2.1) and the bounded-transform theorem give the same cycle. \(\square\)

The numerator is Clifford multiplication by the real position vector. It
is not scalar multiplication by \(|x|\), which would lose both odd
parity and the orientation.

\subsection{3. Differentiation on Euclidean
space}\label{differentiation-on-euclidean-space}

Let \[
H_n=L^2(\mathbb R^n)\otimes\Lambda,\qquad
D_n=\sum_j\widehat c_j\partial_j.
\tag{3.1}
\] Represent \(B_n\) by multiplication by functions and by
\(e_j\mapsto c_j\). Denote this representation by \(\mu\).

\textbf{Lemma 3.1 (the Euclidean Dirac domain).} The closure of (3.1) on
compactly supported smooth forms is odd and self-adjoint, with domain
\(H^1(\mathbb R^n)\otimes\Lambda\), and \(D_n^2=-\Delta\). Its resolvent
is locally compact for \(\mu(B_n)\).

\textbf{Proof.} Under the unitary Fourier transform the operator is
multiplication by the self-adjoint matrix \[
d(\xi)=i\sum_j\xi_j\widehat c_j,\qquad
d(\xi)^2=|\xi|^2.
\] The maximal multiplication domain is exactly the vectors with
\(\int(1+|\xi|^2)\|\widehat u(\xi)\|^2\,d\xi<\infty\). Testing the
adjoint against compactly supported Fourier vectors proves that its
adjoint has the same domain. The matrix inverse resolvents
\((d(\xi)\pm i)^{-1}\) are bounded and onto that domain. Cutoff and
convolution approximation in the first Sobolev norm show that the
compactly supported smooth forms are a core.

For local compactness, first take \(g\in C_c(\mathbb R^n)\). Choose
frequency cutoffs \(h_R\), supported in a ball and tending to one on
expanding balls. The multiplier \[
q_R(\xi)=h_R(\xi)(1+|\xi|^2)^{-1}
\] is square-integrable. The integral kernel of \(M_gq_R(-i\partial)\)
is a constant Fourier-normalization factor times
\(g(x)\check q_R(x-y)\). Its square integral is finite, equal up to that
factor to \(\|g\|_2^2\|q_R\|_2^2\). An \(L^2\) kernel is compact:
approximate it in \(L^2\) by finite sums \(a(x)b(y)\), whose operators
have finite rank, using the kernel \(L^2\) bound on operator norm.

Also \(q_R\to(1+|\xi|^2)^{-1}\) uniformly. Hence \(M_g(1+D_n^2)^{-1}\)
is a norm limit of compact operators. Approximate a general
\(g\in C_0(\mathbb R^n)\) uniformly by compactly supported functions.
The finitely many Clifford matrices are bounded, so the same conclusion
holds for every element of \(B_n\). Adjointing gives the other
localization. \(\square\)

For a compactly supported smooth scalar \(g\), \[
[D_n,M_g]=\sum_j\widehat c_jM_{\partial_jg}.
\tag{3.2}
\] It is bounded, and multiplication by \(g\) preserves the first
Sobolev domain. The graded commutator with each \(c_j\) is zero by
(1.2). The graded Leibniz rule therefore gives bounded commutators for
the dense algebra \(C_c^\infty(\mathbb R^n)\widehat\otimes C_n\).

\textbf{Proposition 3.2 (the Dirac cycle).} The triple \((H_n,\mu,D_n)\)
is an unbounded module, and \[
\alpha_n=[(H_n,\mu,D_n(1-\Delta)^{-1/2})]
\in KK(B_n,\mathbb C).
\]

\textbf{Proof.} Lemma 3.1 proves self-adjointness, its domain and local
compactness. Formula (3.2), the Clifford anticommutation and Sobolev
domain preservation prove the dense smooth commutator condition. Lemma
0.1 now supplies precisely the displayed class. \(\square\)

The individual Dirac resolvent is not globally compact on Euclidean
space. Localizing it in position was essential in Lemma 3.1.

\subsection{4. The first product and its
Gaussian}\label{the-first-product-and-its-gaussian}

The tensor \(B_n\widehat\otimes_\mu H_n\) identifies unitarily with
\(H_n\) by \(b\widehat\otimes u\mapsto\mu(b)u\). The representation is
essential, since function cutoffs converge strongly to one. Under this
identification the first operator becomes \[
X_c=\sum_jx_jc_j.
\] Consider \[
Q_n=D_n+X_c
=\sum_j\bigl((x_j+\partial_j)\varepsilon_j
 +(x_j-\partial_j)\iota_j\bigr).
\tag{4.1}
\] Its domain is the weighted first Sobolev space \[
\mathcal W=\left\{u\in H_n:
x_ju,\ \partial_ju\in H_n\ \text{for every }j\right\}.
\tag{4.2}
\]

\textbf{Lemma 4.0 (the complete oscillator input).} The closure of
\(Q_n\) on compactly supported smooth forms is self-adjoint on precisely
\(\mathcal W\). Its inverse resolvents are compact. Its kernel is the
even line generated by \(e^{-|x|^2/2}\); its even-to-odd part is onto.

\textbf{Proof.} Cut off and mollify to approximate any vector of
\(\mathcal W\) in all the norms \(\|u\|,\|x_ju\|,\|\partial_ju\|\). For
convolution use
\(x_j(\rho_\epsilon*u)=\rho_\epsilon*(x_ju)+(x_j\rho_\epsilon)*u\),
whose last term has norm at most \(C\epsilon\|u\|\). Thus the stated
smooth space is dense in \(\mathcal W\).

On that core the exterior and Clifford relations give, with
\(a_j=\partial_j+x_j\), \[
\|Q_nu\|^2=\sum_j\|a_ju\|^2+2\langle\mathcal Nu,u\rangle
=\sum_j(\|x_ju\|^2+\|\partial_ju\|^2)
 +\langle(2\mathcal N-n)u,u\rangle.
\tag{4.0a}
\] Integration by parts gives the cross term
\(2\operatorname{Re}\langle x_ju,\partial_ju\rangle=-\|u\|^2\). Since
\(0\leq\mathcal N\leq n\), the graph norm and the weighted Sobolev norm
are equivalent. Completeness and density identify the closure domain
with \(\mathcal W\), and extend (4.0a) to it.

Here is self-adjointness without an elliptic theorem left implicit. If
\(Q_n^*u=\pm iu\), for a smooth compactly supported cutoff \(\chi\) we
have in distributions
\(D_n(\chi u)=\chi(\pm iu-X_cu)+[D_n,\chi]u\in L^2\). The Fourier matrix
in Lemma 3.1 then implies \(\chi u\in H^1\). In particular
\(\chi u\in\mathcal W\). Choose \(\chi_R(x)=\chi(x/R)\), equal to one on
the radius-\(R\) ball, with \(\|[D_n,\chi_R]\|\leq C/R\). Symmetry on
\(\mathcal W\) and \(Q_n(\chi_Ru)=\pm i\chi_Ru+[D_n,\chi_R]u\) give
\(\|\chi_Ru\|^2\leq(C/R)\|\chi_Ru\|\|u\|\). Let \(R\to\infty\); hence
\(u=0\). The ranges of \(Q_n\pm i\) are closed, since their norms
squared are \(\|Q_nu\|^2+\|u\|^2\); their orthogonal complements are
these zero deficiency kernels. Both are onto, proving self-adjointness.

The inclusion \(\mathcal W\to H_n\) is compact. Indeed choose bounded
position and frequency cutoffs \(\chi_R,h_R\), equal to one on
radius-\(R\) balls and zero outside radius-\(2R\) balls. The operator
\(M_{\chi_R}h_R(-i\partial)\) has an \(L^2\) kernel and is compact by
the finite-rank kernel proof in Lemma 3.1. Its error on the unit ball of
\(\mathcal W\) is at most \[
\|(1-\chi_R)u\|+\|\chi_R(1-h_R(-i\partial))u\|
\leq R^{-1}(\||x|u\|+\||\partial|u\|).
\tag{4.0b}
\] Thus the inclusion is a norm limit of compact maps. Inverse
resolvents map boundedly into the graph domain, so are compact.

If \(Q_nu=0\), the first expression in (4.0a) forces \(\mathcal Nu=0\)
and \(a_ju=0\) for every \(j\). The first says that only the degree-zero
coefficient remains. Locally the latter equations say that the
distribution \(e^{|x|^2/2}u\) has all derivatives zero, and hence is
constant: convolve on a smaller ball, use the elementary zero-gradient
assertion there, and pass to distributions; constants agree on
overlapping balls. The resulting Gaussian lies in \(\mathcal W\) and is
even. Finally \(Q_n\) is bounded below on its kernel's orthogonal
complement. Otherwise unit vectors there with \(Q_nu_m\to0\) would be
bounded in \(\mathcal W\), have a norm-convergent subsequence by
compactness, and limit to a unit vector in the kernel and its orthogonal
complement. Its range is therefore closed and equals the kernel's
orthogonal complement by self-adjointness. Since the kernel is entirely
even, the even-to-odd block is onto. \(\square\)

\textbf{Theorem 4.1 (the first inverse identity).} One has \[
\beta_n\widehat\otimes_{B_n}\alpha_n=1_{\mathbb C}.
\tag{4.3}
\]

\textbf{Proof.} Compactness of the graph inclusion in (4.2) makes the
inverse resolvents of \(Q_n\) compact: they map boundedly into that
graph domain. Thus \(Q_n\), with scalar source, is an unbounded cycle.

Check the product criterion on the dense creation vectors
\(b\in C_c^\infty(\mathbb R^n)\widehat\otimes C_n\). Their creation maps
are \(\mu(b)\). The graded connection error is \[
Q_n\mu(b)-(-1)^{|b|}\mu(b)D_n
=[D_n,\mu(b)]_{\rm gr}+X_c\mu(b).
\tag{4.4}
\] It is bounded: the first term is (3.2) with Clifford factors, and the
second is a compactly supported matrix function. Multiplication by such
\(b\) sends the first Sobolev domain into (4.2); the adjoint sends (4.2)
into the first Sobolev domain. This proves both required domain
conditions.

The domain in (4.2) is contained in \(\operatorname{Dom}X_c\), and \[
\{D_n,X_c\}=\sum_j\widehat c_jc_j=2\mathcal N-n,
\qquad
\mathcal N=\sum_j\varepsilon_j\iota_j .
\tag{4.5}
\] Indeed the derivative terms containing \(x_k\partial_j\) cancel by
(1.2); differentiating \(x_j\) gives the first sum, and
\((\varepsilon_j-\iota_j)(\varepsilon_j+\iota_j)=2\varepsilon_j\iota_j-1\).
Consequently the positivity form is \[
2\operatorname{Re}\langle X_cu,Q_nu\rangle
=2\||x|u\|^2+\langle u,(2\mathcal N-n)u\rangle
\geq-n\|u\|^2 .
\tag{4.6}
\] The form identity, initially on compactly supported smooth forms,
extends to (4.2) by its graph norm. Theorem 0.5 identifies the bounded
transform of \(Q_n\) with the product.

By Lemma 4.0 the kernel is the even Gaussian line. On its orthogonal
complement replace the bounded transform by its sign. This is a compact
change: there is a spectral gap at zero by the last argument of that
lemma, and the difference is a continuous function on the spectrum
vanishing at infinity. The compact resolvent makes every such function
compact, by the cutoff-functional-calculus argument in Lemma 0.1. The
complement is now an odd self-adjoint involution with scalar
representation and hence a degenerate cycle. The remaining even
one-dimensional zero-operator cycle is precisely \(1_{\mathbb C}\). In
particular its even-to-odd Fredholm index is \(+1\), agreeing with the
earlier scalar calculation, but no scalar-picture isomorphism is needed
to prove (4.3). \(\square\)

For \(n=1\) the even-to-odd operator is \(x+d/dx\). Its equation says
\((e^{x^2/2}u)'=0\), giving the Gaussian. The formal adjoint \(x-d/dx\)
instead has solutions proportional to \(e^{x^2/2}\), none in \(L^2\)
except zero. Lemma 4.0 proves the necessary closed-range and domain
assertions in addition to these differential equations.

\subsection{5. Rotating the two position
spaces}\label{rotating-the-two-position-spaces}

The first identity alone does not imply the second. We prove the missing
identity by a homotopy of cycles, rather than assuming cancellation.

Write \[
B_n\widehat\otimes B_n
=C_0(\mathbb R_x^n\times\mathbb R_y^n)
\widehat\otimes C_{2n},
\] with generators \(e_j\) in the first factor and \(f_j\) in the
second. They anticommute with each other. Define the even graded
automorphism \(\rho\) of \(B_n\) by \[
\rho(g)(x)=g(-x),\qquad \rho(e_j)=-e_j.
\tag{5.1}
\] Both the coordinates and Clifford generators are reflected.

\textbf{Lemma 5.1 (the rotation relation).} In
\(KK(B_n,B_n\widehat\otimes B_n)\), \[
\beta_n\widehat\otimes1_{B_n}
=[\rho]\widehat\otimes\beta_n .
\tag{5.2}
\]

\textbf{Proof.} Let \(0\leq\theta\leq\pi/2\), and write
\(c=\cos\theta\), \(s=\sin\theta\). Rotate coordinates and generators
together: \[
\begin{gathered}
u=cx+sy,\qquad v=-sx+cy,\\
e_j^\theta=ce_j+sf_j,\qquad
f_j^\theta=-se_j+cf_j.
\end{gathered}
\tag{5.3}
\] The new generators satisfy the same Clifford relations; the two
families anticommute. On the standard module \(B_n\widehat\otimes B_n\),
represent the source \(B_n\) by \(g(v)\) and \(e_j\mapsto f_j^\theta\),
and take \[
F_\theta(x,y)=
\frac{\sum_ju_je_j^\theta}{\sqrt{1+|u|^2}}.
\tag{5.4}
\] Its source graded commutators are zero, since the Clifford families
are perpendicular, and \[
(F_\theta^2-1)\phi_\theta(a)
=-\frac{\phi_\theta(a)}{1+|u|^2}
\] is compact. To check the homotopy, use the interval standard module
with these jointly continuous fields. Its multipliers are strictly
continuous and uniformly bounded. The localized defects belong to
\(C([0,\pi/2],B_n\widehat\otimes B_n)\): when \(|(x,y)|\) is large,
either \(|u|\) is large, making the resolvent factor small, or \(|v|\)
is large, making the source function small. This estimate is uniform in
\(\theta\), since the rotation preserves \(|x|^2+|y|^2\). Approximation
by finite Clifford sums proves it for every source element.

Thus these fields form a genuine interval Kasparov cycle. At zero the
source is in the second factor and the operator is the first Bott
multiplier, representing \(\beta_n\widehat\otimes1\). At \(\pi/2\) the
source is \(g(-x)\) with generators \(-e_j\), and the operator is the
second Bott multiplier. This is \([\rho]\widehat\otimes\beta_n\).

The fields need not form a path continuous in operator norm on a fixed
module: the rotating coordinate is unbounded. The interval multiplier
and uniform localized-defect checks above are the appropriate homotopy
assertions. \(\square\)

\textbf{Theorem 5.2 (the second inverse identity).} One has \[
\alpha_n\widehat\otimes_{\mathbb C}\beta_n=1_{B_n}.
\tag{5.5}
\]

\textbf{Proof.} Compose (5.2) on the right with
\(1_{B_n}\widehat\otimes\alpha_n\). The interchange rule for exterior
products and associativity identify the left side with
\(\alpha_n\widehat\otimes_{\mathbb C}\beta_n\): the second-factor input
is evaluated by \(\alpha_n\), and the first-factor output is produced by
\(\beta_n\). The right side becomes
\([\rho]\widehat\otimes(\beta_n\widehat\otimes_{B_n}\alpha_n)=[\rho]\),
by (4.3). Hence \(p=\alpha_n\widehat\otimes_{\mathbb C}\beta_n=[\rho]\).

Associativity and (4.3) also give \[
p\widehat\otimes_{B_n}p
=\alpha_n\widehat\otimes
(\beta_n\widehat\otimes_{B_n}\alpha_n)
\widehat\otimes\beta_n=p .
\] But \([\rho]\) is invertible: \(\rho^2=1\). An invertible idempotent
in a unital ring is the identity, by multiplying \(p^2=p\) by its
inverse. Thus \(p=1_{B_n}\). This proves the second identity and also
shows \([\rho]=1_{B_n}\) in KK. \(\square\)

This proof uses the graded Bott algebra. Reflecting the coordinate alone
is a different automorphism; no assertion about its class was made in
the rotation.

\subsection{6. The planar projection and its
sign}\label{the-planar-projection-and-its-sign}

Represent \(C_2\) on \(\mathbb C^2\) by \(e_1=\sigma_1\),
\(e_2=\sigma_2\), with grading \(\sigma_3\). Graded Morita equivalence
identifies \(B_2\) with \(C_0(\mathbb R^2)\). In the oriented coordinate
\(z=x_1+ix_2\), the Bott transform is \[
F_2(z)=\frac1{\sqrt{1+|z|^2}}
\begin{pmatrix}0&\overline z\\z&0\end{pmatrix}.
\tag{6.1}
\] Its even-to-odd corner is \(T=z/\sqrt{1+|z|^2}\).

\textbf{Proposition 6.1 (the positive planar Bott class).} Under the
scalar KK-to-\(K_0\) index identification, (6.1) gives \[
\mathfrak b=[q]-[P],\qquad
q(z)=\frac1{1+|z|^2}
\begin{pmatrix}1&\overline z\\z&|z|^2\end{pmatrix},
\quad P=\begin{pmatrix}0&0\\0&1\end{pmatrix}.
\tag{6.2}
\] This is the Bott projection convention of \emph{Vector bundles and
finitely generated projective modules} and the positive
double-suspension convention of \emph{Bott periodicity in operator
K-theory}.

\textbf{Proof.} Put \(r=(1+|z|^2)^{-1/2}\). Since \(T^*T=TT^*=1-r^2\),
the multiplier \[
W=\begin{pmatrix}T&-r\\r&T^*\end{pmatrix}
\] is unitary. Its first-column projection is \[
W\begin{pmatrix}1&0\\0&0\end{pmatrix}W^*
=\frac1{1+|z|^2}
\begin{pmatrix}|z|^2&z\\\overline z&1\end{pmatrix}.
\tag{6.3}
\] The scalar Fredholm-picture index formula is this projection minus
its scalar projection at infinity. It is the same formula that gives
\([1-SS^*]\) with a negative sign for a unilateral shift \(S\): the
Julia unitary lift is
\(\left(\begin{smallmatrix}S&-(1-SS^*)\\0&S^*\end{smallmatrix}\right)\),
whose first-column projection is \(\operatorname{diag}(SS^*,0)\). Thus
its difference from \(\operatorname{diag}(1,0)\) is \(-[1-SS^*]\). This
fixes the index-boundary sign. Here the cycle-to-relative-projection
identification can be checked directly. If \(R\) is (6.3) and
\(P_\infty=\operatorname{diag}(1,0)\), use the cycle on
\(B^2\oplus B^2\) with source projections \(R,P_\infty\) and
off-diagonal identity. Their difference is compact, so it is a cycle.
Removing its zero-source parts changes its operator only compactly; the
remaining off-diagonal map is \(P_\infty:R B^2\to P_\infty B^2\). The
unitary column \(b\mapsto(Tb,rb)\) identifies that map with \(T\). Thus
the relative projection pair gives exactly (6.1), including the
prescribed boundary sign, without using the unfinished Fredholm-picture
lesson.

Conjugate (6.3) and its scalar projection by the constant
vector-coordinate swap. The resulting pair is exactly \(q,P\) in (6.2),
so the relative class is \([q]-[P]\).

For an independent check of the orientation, its finite-plane line frame
is \((1,z)^T\), and its infinity-disc frame is \((1/z,1)^T\). The first
equals \(z\) times the second. Thus its coefficient transfer on the
equator is multiplication by \(z\), of winding \(+1\), exactly the
convention in the vector-bundle lesson. Section 6 of the operator
K-theory periodicity lesson identifies this relative class with the
positive ordered double suspension. Complex conjugating \(z\), or
interchanging the two real coordinates, reverses this winding and the
class. Swapping the two vector coordinates as above preserves the class,
because the scalar projection is swapped with it. \(\square\)

\subsection{7. The one-dimensional extension
cycles}\label{the-one-dimensional-extension-cycles}

Put \(S=C_0(\mathbb R)\), and identify the real coordinate with
\(0<t<1\) by \[
t(s)=\frac12+\frac1\pi\arctan s,\qquad
u(s)=e^{2\pi it(s)}=\frac{s-i}{s+i}.
\tag{7.1}
\] The unitary \(u\) winds positively once as \(s\) increases. We
specify the left Clifford transfer before comparing extensions.

A right \(C_1\)-cycle, regarded as a module over its ungraded
coefficient algebra, has the odd right multiplication \(R_\varepsilon\)
and grading \(\Gamma\). The odd self-adjoint unitary \[
J=iR_\varepsilon\Gamma
\tag{7.2}
\] is the left Clifford action: it anticommutes with the cycle operator,
because that operator commutes with \(R_\varepsilon\) and anticommutes
with \(\Gamma\). The coefficient inner product here is obtained by
taking the even part of the \(C_1\)-valued inner product. For the
inverse construction first take the unital essential source summand. Its
orthogonal complement has zero source; compact off-diagonal errors can
be removed by the cycle perturbation rule. On that summand replace the
operator by \((F-JFJ)/2\), a compact perturbation making its
anticommutation with \(J\) exact. The inverse construction now recovers
right multiplication as \(-iJ\Gamma\), which commutes with this
normalized operator, and recovers the coefficient inner product as \[
\langle\xi,\eta\rangle_S+
\langle\xi,R_\varepsilon\eta\rangle_S\varepsilon.
\] Positivity follows by evaluating \(\varepsilon\) at \(+1\) and
\(-1\): the two values are the positive forms of \(1\pm R_\varepsilon\).
The norms before and after taking the even part are equivalent, so
completions agree.

Compactness is preserved in both directions. A coefficient rank-one
operator becomes the sum
\(\theta_{\xi,\eta}^S+R_\varepsilon\theta_{\xi,\eta}^SR_\varepsilon\).
Conversely a compact \(S\)-operator commuting with \(R_\varepsilon\) is
approximated by such sums: average its finite-rank approximants with
their \(R_\varepsilon\)-conjugates. These observations apply to defects,
commutators, and homotopies. Thus (7.2) gives the required formal
Clifford-transfer isomorphism.

In the standard right module \(S\widehat\otimes C_1\),
\(R_\varepsilon=\sigma_1\), \(\Gamma=\sigma_3\), and \(J=\sigma_2\).
Conjugate by the even unitary \(\exp(i\pi\sigma_3/4)\). The left
generator becomes \(\sigma_1\), while the operator \(h(s)R_\varepsilon\)
becomes \(-h(s)\sigma_2\). This minus sign is part of the transfer
convention.

\textbf{Proposition 7.1 (the cone element).} The extension \[
0\longrightarrow S\longrightarrow C_0((0,1])
\xrightarrow{\operatorname{ev}_1}\mathbb C\longrightarrow0
\tag{7.3}
\] gives the element \(x\in KK(C_1,S)\) represented, after (7.2), by \[
\left(S\otimes\mathbb C^2,\,
\varepsilon\mapsto\sigma_1,\,
-\frac{s}{\sqrt{1+s^2}}\sigma_2\right),
\quad\Gamma=\sigma_3.
\tag{7.4}
\]

\textbf{Proof.} We identify its extension cycle directly. On \(S^2\) put
\(Vb=(\sqrt t\,b,\sqrt{1-t}\,b)\), \(P=VV^*\), and represent the scalar
quotient by \(\rho(\lambda)=\operatorname{diag}(\lambda,0)\). This is
the Stinespring dilation of the completely positive multiplier lift
\(L(\lambda)=t\lambda\): \(V^*\rho(\lambda)V=t\lambda\). The operator
\(2P-1\) is an involution and its source commutator is compact, because
its off-diagonal entries are \(2\sqrt{t(1-t)}\in S\). Tensoring with the
right odd generator gives the extension cycle. This dilation realizes
the actual extension: the Busby image of the positive lift \(t\) is a
projection, and its lift is precisely \(L(1)\).

Remove the compact off-diagonal entries of \(2P-1\). The first diagonal
block is \(2t-1\), and the second has zero source representation and is
degenerate. Compact perturbation and deletion of the zero-source cycle
therefore give the ordinary odd cycle \((S,\lambda,2t-1)\). The
functions \(2t(s)-1=(2/\pi)\arctan s\) and \(h(s)=s/\sqrt{1+s^2}\) have
the same limits \(-1,+1\). Their straight interpolation is a
norm-continuous path of self-adjoint contraction multipliers with square
defect in \(S\). Applying the explicitly proved Clifford transfer (7.2)
gives (7.4).

For orientation, the positive exponential of the lift is
\(e^{2\pi it}=u\). Formula (7.1) computes its positive winding directly.
Thus the calculation specifies the cone class and its coordinate.
\(\square\)

Let \(P=-i\,d/ds\) on \(L^2(\mathbb R)\), and let \(H_-\) be its
negative-frequency subspace.

\textbf{Proposition 7.2 (the coisometry extension).} The odd cycle with
ungraded operator \(-P(1+P^2)^{-1/2}\) gives the extension \[
0\longrightarrow\mathcal K\longrightarrow C^*(v-1)
\longrightarrow S\longrightarrow0,
\tag{7.5}
\] where \(v\) is a coisometry of index \(+1\). Its Busby map satisfies
\(\tau(u-1)=q(v)-1\). Denote its class in \(KK(S,C_1)\) by \(y\).

\textbf{Proof.} Replace its bounded transform locally compactly by the
sign of \(-P\). Although the sign is discontinuous at zero, this is a
valid locally compact change: its difference from the bounded transform
is \(g(P)(1+P^2)^{-1}\) for a bounded Borel function \(g\), since the
difference decays quadratically at infinity. Lemma 3.1 makes the inverse
resolvent locally compact; multiplication by the bounded Borel operator
preserves that compactness. The positive spectral projection of the sign
is precisely the projection onto \(H_-\). The compression
\(L(a)=P_-M_aP_-\) on that subspace is completely positive and
contractive; its multiplicative errors are compact because \([P_-,M_a]\)
is compact. The original multiplication representation on
\(L^2(\mathbb R)\), with \(V\) the inclusion of \(H_-\), is a dilation
of \(L\). Its operator \(2VV^*-1\) is exactly the sign just obtained.
Hence this is the extension cycle of its compression Busby map.
Independence of the choice of dilation follows by the explicit rotation
of two isometries with common compression:
\(V_\theta=(\cos\theta V_0,\sin\theta V_1)\); the compact errors
\(\rho_i(a)V_i-V_iL(a)\) make the projection commutators compact, and
the unused endpoint summand is degenerate. This is a local verification,
not an appeal to a picture theorem.

Here is an explicit basis identifying the compression. On the circle
with normalized measure define the unitary \[
(Uf)(s)=\frac{1}{\sqrt\pi(s-i)}f(u(s)):
L^2(\mathbb T)\longrightarrow L^2(\mathbb R).
\tag{7.6}
\] Indeed \(d\theta/ds=2/(1+s^2)\), so the change of variables gives
equality of the two squared norms, and the formula has a measurable
inverse.

The images \(Uz^{-m}\), \(m\geq0\), have negative Fourier support, while
\(Uz^k\), \(k\geq1\), have positive support. To verify this directly,
the first functions are proper rational functions with their only pole
at \(i\), and the second are proper rational functions with their only
pole at \(-i\). Expand each in powers of the inverse linear factor. The
identity \[
\frac1{s-i}=i\int_0^\infty e^{-r}e^{-isr}\,dr
\] and its derivatives give negative Fourier support for every inverse
power; the conjugate identity gives positive support at the other pole.
These formulas hold in \(L^2\), by the Fourier transform of the
integrable exponential-polynomial functions. Since all the \(Uz^k\),
\(k\in\mathbb Z\), form a complete orthonormal basis, the two stated
spans are exactly \(H_-\) and its orthogonal positive-frequency
complement.

In the basis \(Uz^{-m}\) of \(H_-\), compression of \(M_u\) sends the
\(m\)-th vector to the \((m-1)\)-st for \(m\geq1\), and sends the zeroth
to zero. It is the adjoint unilateral shift \(v\). Thus \(vv^*=1\), its
kernel is one dimensional, and its index is \(+1\). The algebra
generated by \(v-1\) contains the compact operators: its elements
satisfy \[
(v^*-1)(v-1)+(v-1)+(v^*-1)=-(1-v^*v).
\] Thus it contains the kernel projection, and multiplying this
projection by \(1+(v-1)\) and \(1+(v^*-1)\) gives all finite matrix
units without requiring the unit to belong to the algebra. Its quotient
is the functions on the circle vanishing at \(1\), namely \(S\); the
shift Calkin spectrum is the whole circle. To verify this last fact
here, a unit vector with coefficients \(1,z,\ldots,z^{N-1}\) on a block
of \(N\) consecutive basis vectors, normalized by \(N^{-1/2}\), has
\(\|(v-z)\xi\|\leq 2/\sqrt N\) for \(|z|=1\). Take disjoint blocks
escaping to infinity; the vectors are weakly zero, so compact operators
send them to zero. A Calkin inverse of \(v-z\) would contradict this
estimate. The quotient image of \(v\) is unitary, so has no other
spectrum. Continuous functional calculus then gives the full circle
quotient. This proves (7.5) and the asserted Busby map. The displayed
minus one is subtraction of the unit. \(\square\)

\textbf{Theorem 7.3 (the two extension products).} In these conventions,
\[
x\widehat\otimes_S y=1_{C_1},\qquad
y\widehat\otimes_{C_1}x=1_S.
\tag{7.7}
\]

\textbf{Proof.} Transport the inverse pair \(\beta_1,\alpha_1\) by
exterior tensoring with \(C_1\) and the fixed \(C_2\) Morita
equivalence. Let \(\widetilde x\in KK(C_1,S)\) be the resulting Bott
element. Its left generator is \(\sigma_1\), its grading \(\sigma_3\),
and its operator is \(+h\sigma_2\): in the ordered tensor
\(C_1^{\rm aux}\widehat\otimes(S\widehat\otimes C_1^{\rm orig})\),
represent the auxiliary generator by \(\sigma_1\) and the original one
by \(\sigma_2\). The exterior Bott operator is \(h\) times the latter,
and the source acts through the former. Its inverse
\(\widetilde y\in KK(S,C_1)\) comes from the opposite tensor order in
the inverse construction. We compute that order to fix its sign.

On \(L^2(\mathbb R)\otimes\mathbb C^2\widehat\otimes C_1\), its two left
Clifford generators, in order auxiliary then original, are \[
\gamma_1=\sigma_3L_\varepsilon,\qquad
\gamma_2=\sigma_1,\qquad
\chi=-i\gamma_1\gamma_2=\sigma_2L_\varepsilon.
\] The operator is \(\sigma_2P\). The even Morita corner is
\((1+\chi)/2\). As a right \(C_1\)-module its range has the homogeneous
even unit vector \[
\xi_0=\frac1{\sqrt2}
\begin{pmatrix}1\\ i\varepsilon\end{pmatrix}.
\] Its rank-one projection is exactly \((1+\chi)/2\), and
\(\sigma_2P\xi_0=\xi_0\varepsilon P\). Therefore \(\widetilde y\) is the
right odd cycle with ungraded operator \(+P\). The two tensor-unitary
Morita identities and associativity give
\(\widetilde x\widetilde y=1_{C_1}\) and
\(\widetilde y\widetilde x=1_S\), directly from (4.3) and (5.5).

By (7.4), \(x=-\widetilde x\). Explicitly, the additive inverse is
obtained by reversing the grading; the odd unitary \(\sigma_1\) carries
that inverse cycle back to grading \(\sigma_3\), keeps the source
generator fixed, and changes \(h\sigma_2\) to \(-h\sigma_2\).
Proposition 7.2 similarly gives \(y=-\widetilde y\) in the right odd
picture. Bilinearity cancels the two signs, proving both identities in
(7.7). This proves the extension assertions with their actual cone,
coisometry and Clifford conventions. \(\square\)

\subsection{8. Periodicity from the inverse
cycles}\label{periodicity-from-the-inverse-cycles}

For separable graded \(A\) and separable coefficient algebra \(B\),
exterior product with the two inverse cycles gives inverse maps \[
\begin{aligned}
KK(A,B)&\longrightarrow KK(A,B\widehat\otimes B_n),\\
x&\longmapsto x\widehat\otimes\beta_n,\\
KK(A,B\widehat\otimes B_n)&\longrightarrow KK(A,B),\\
y&\longmapsto y\widehat\otimes_{B_n}\alpha_n,
\end{aligned}
\tag{8.1}
\] where the second map includes the identity on \(B\). The two
composites are identity by (4.3), (5.5), associativity and the exterior
interchange rule. Their analogues in the first variable are likewise
inverse.

Since \(B_n=C_0(\mathbb R^n)\widehat\otimes C_n\), putting \(n=1\)
trades a suspension for one Clifford degree. Putting \(n=2\), and using
the graded Morita equivalence \(C_2\sim\mathbb C\), gives the natural
two-periodicity isomorphisms \[
KK(A,B)\cong KK(A,S^2B)\cong KK(S^2A,B).
\tag{8.2}
\] For \(n=1\) it gives the corresponding natural degree-one
isomorphisms \[
KK^1(A,B)\cong KK(A,SB)\cong KK(SA,B).
\tag{8.3}
\] The maps are products with fixed cycles, not merely abstract group
identifications. Therefore homomorphism naturality follows from the
product naturality theorem.

\textbf{Proposition 8.1 (arbitrary coefficient algebras).} For separable
\(A\), the isomorphisms (8.1)--(8.3) hold for every coefficient algebra
\(B\), without a separability or \(\sigma\)-unitality assumption on
\(B\).

\textbf{Proof.} We prove the coefficient reduction used here. Given a
cycle on a countably generated Hilbert \(B\)-module \(E\), choose a
countable homogeneous generating family and a countable dense
homogeneous family in \(A\). Enlarge the vector family successively by
the images under the cycle operator, its adjoint, the grading, and the
represented source family and its adjoints. Include also all the vectors
in finite rank-one approximations, to arbitrary accuracy, of the cycle
defects for the countably many source elements. Repeat these operations
countably many times.

Let \(B_0\subset B\) be the separable graded C*-subalgebra generated by
all inner products of the resulting countable family. Let \(E_0\) be the
closure in \(E\) of its \(B_0\)-span. The inner products take values in
\(B_0\). Every selected vector belongs to \(E_0\): a countable
approximate identity of \(B_0\) approximates it in norm, since its own
squared inner product belongs to \(B_0\). Thus \(E_0\) is countably
generated over \(B_0\).

The operator, its adjoint, grading and dense source family preserve
\(E_0\), by construction. Continuity extends the source action to all
\(A\). The chosen finite rank-one approximations make the restricted
defects compact on \(E_0\); continuity gives the cycle axioms for every
source element. The inner-product formula makes \[
E_0\widehat\otimes_{B_0}B\longrightarrow E,\qquad
\xi\widehat\otimes b\longmapsto\xi b
\tag{8.4}
\] isometric. Its range contains the original generating vectors times
\(B\), so it is onto. It intertwines the source and the operator. Every
class therefore comes from a separable coefficient subalgebra.

The same reduction applies to a homotopy module over \(C([0,1],B)\). All
its selected inner products are continuous \(B\)-valued functions. The
values of this countable family at a countable dense set of parameters
generate a separable subalgebra \(B_0\); continuity puts every value in
\(B_0\). The construction with right coefficients \(C([0,1],B_0)\) then
descends the entire homotopy. To preserve given endpoint descents, lift
their countable generating vectors to the homotopy module: evaluation
onto a fibre is the quotient by the submodule generated by functions
vanishing at that endpoint, so every fibre vector has a lift. Include
those lifts in the selection. Also include the original endpoint
coefficient subalgebras and the images of those vectors under the
endpoint intertwining unitaries. At a later common separable coefficient
stage, the endpoint tensor maps are therefore the prescribed unitaries,
by their inner-product formulas and density. Finite sums and inverse
cycles descend as well. Consequently \[
KK(A,B)=\underset{B_0\subset B\ {\rm separable}}{\operatorname{colim}}KK(A,B_0).
\tag{8.5}
\] Here separable graded subalgebras form a directed system: the algebra
generated by two countable dense subsets and their grading images is
again separable.

The periodicity maps already proved for separable \(B_0\) commute with
coefficient homomorphisms. Pass to (8.5). Every separable subalgebra of
\(S^nB\) is contained in \(S^nB_0\) for some separable \(B_0\), by
applying the same countable-values argument to a dense family of its
functions. Hence the target colimits identify with the required
suspended coefficient groups, and the two inverse identities survive in
the colimit. For first-variable suspension the source \(S^nA\) is still
separable, so the identical argument applies. This proves the stated
arbitrary-coefficient generality. \(\square\)

\subsection{9. Exercises}\label{exercises}

\textbf{9.1. The Bott defect.} Verify every bounded-cycle condition for
(2.2), including countable generation of its module and compactness of
the square defect.

\textbf{9.2. The one-dimensional oscillator.} Compute the kernels of
\(x+d/dx\) and its formal adjoint on the weighted first Sobolev domain,
and justify why the index is \(+1\).

\textbf{9.3. The rotation.} Give the interval module, source
representation and operator in the rotation relation. Explain the
uniform compactness estimate and why the reflected Clifford action must
accompany the coordinate reflection.

\textbf{9.4. The planar sign.} Compute the relative projection of the
even-to-odd corner in (6.1), compare it with the fixed vector-bundle
Bott projection, and determine the effect of reversing the orientation
of the real plane.

\subsection{10. Solutions}\label{solutions}

\textbf{Solution to 9.1.} The Clifford relations give
\((\sum_jx_je_j)^2=|x|^2\), so the odd self-adjoint multiplier is a
contraction and has defect \(-(1+|x|^2)^{-1}\). That function belongs to
\(B_n\), and left multiplication by it is compact on the standard
module, whose rank-one operators are multiplication by \(bd^*\). Scalar
commutators and the adjoint defect are zero. A countable approximate
identity of the separable algebra \(B_n\) generates the module. Thus
every Kasparov-cycle condition holds.

\textbf{Solution to 9.2.} The distributional equations give respectively
\(u=Ce^{-x^2/2}\) and \(u=Ce^{x^2/2}\), by multiplication by the
relevant integrating factor. The first lies in the weighted domain; the
second lies in \(L^2\) only for \(C=0\). Lemma 4.0 proves that
\(x+d/dx\) has closed and surjective range on that domain. Equivalently
its graph estimate and compact graph inclusion give closed range, while
its adjoint-kernel computation gives dense range. Its kernel is one
dimensional and its cokernel zero. The index is \(+1\).

\textbf{Solution to 9.3.} Use \(C([0,\pi/2],B_n\widehat\otimes B_n)\),
the source \(g(v),e_j\mapsto f_j^\theta\), and (5.4). The generators are
orthogonal after the simultaneous rotation, so all graded source
commutators vanish. The square defect times a source element decays
uniformly: \(|u|^2+|v|^2=|x|^2+|y|^2\), so outside a sufficiently large
ball either the inverse quadratic factor is small or the source function
is small. Joint continuity and strict multiplier continuity give the
interval cycle. Its endpoints are those in (5.2). At the terminal angle
both \(v=-x\) and \(f_j^\theta=-e_j\); retaining only the coordinate
reflection would describe a different endpoint. Compose with
\(1\widehat\otimes\alpha_n\), then use the invertible-idempotent
argument of Theorem 5.2 to obtain the reverse product.

\textbf{Solution to 9.4.} For \(T=z/\sqrt{1+|z|^2}\) and
\(r=(1+|z|^2)^{-1/2}\), the unitary \(W\) in Section 6 has first-column
projection (6.3). Subtract its value \(\operatorname{diag}(1,0)\) at
infinity, then conjugate both projections by the constant coordinate
swap. This gives precisely \([q]-[P]\) of (6.2). Its frames satisfy
\((1,z)=z(1/z,1)\); the coefficient-transfer winding is \(+1\).
Replacing \(z\) by \(\overline z\) reverses it to \(-1\), so reverses
the relative class. The operator K-theory lesson's ordered suspension
has coordinate \(x+iy\), so the sign in (6.2) agrees with that positive
generator.

\subsection{Working completion
obligations}\label{working-completion-obligations}

The Bott and Dirac cycles, both inverse products, the planar projection
sign, and the exact cone/coisometry extension comparison are given
above. Proposition 8.1 also proves arbitrary-coefficient periodicity for
separable sources. Periodicity for nonseparable sources remains to be
checked against the exact homotopy and product conventions used earlier
in this course. The scalar cone/coisometry comparison is proved above.
The full-scope nonseparable-source periodicity assertion and its uses
still require the indicated further proof.

\subsection{References and source
credit}\label{references-and-source-credit}

\begin{itemize}
\tightlist
\item
  B. Blackadar, \emph{K-Theory for Operator Algebras},
  \href{https://www.bruceblackadar.com/Mathematics/book6.pdf}{actual
  free corrected author edition}, Sections 19.2 and 19.9.3, provides the
  scalar cone/coisometry conventions and the Euclidean-cycle comparison.
  The inverse products, weighted domain, Gaussian and Cayley compression
  are proved above. The
  \href{https://www.bruceblackadar.com/mathpubs.html}{author's
  publications page} retains the author's copyright and permits use
  under Creative Commons rules with attribution, without specifying a
  particular licence or version; no more specific licence is asserted
  here.
\item
  K. van den Dungen,
  \href{https://arxiv.org/pdf/2006.10616v2}{\emph{Localisations of
  half-closed modules and the unbounded Kasparov product}, actual arXiv
  version 2006.10616v2}, Sections 2.4 and 3.1, printed pp.~11--16, is
  the free comparison source for Section 0. The domain, connection,
  integrated-positivity and nonessential-representation arguments used
  here are given in full locally. Availability of this preprint does not
  assert a separate open adaptation licence.
\item
  G. Teschl, \emph{Mathematical Methods in Quantum Mechanics},
  \href{https://www.mat.univie.ac.at/~gerald/ftp/book-schroe/schroe.pdf}{actual
  free author-hosted online edition}, dated 12 February 2009, Section
  7.1, Lemmas 7.1 and 7.3 and Theorems 7.4--7.5, printed pp.~161--164,
  is the comparison account for the locally proved Euclidean Fourier
  argument. The file retains © 2009 American Mathematical Society and
  states that this online edition is available for personal online use;
  it is free to read, with no open republication or adaptation licence
  asserted here.
\end{itemize}

\end{document}
